\section*{الفصل \ref{ch:b3:glaciology} --- علم الجليد}

\begin{solution}{exo:b3:glaciology:1}
$\tau_b = \rho gH\sin\alpha = 917\times9.81\times250\times\sin\ang{3} =
\qty{118}{kPa}$، أي حوالي \qty{1.2}{bar}: وهو من رتبة الإجهادات تحت
الأنهار الجليدية الحقيقية.
\end{solution}

\begin{solution}{exo:b3:glaciology:2}
(أ)~$\qty{2.9e15}{m^3}\times917 = \qty{2.66e18}{kg} = \qty{2.66e6}{Gt}$؛
مقسومةً على \qty{360}{Gt/mm}، تعطي \qty{7390}{mm}، أي \qty{7.4}{m}، القيمةَ المذكورة.
(ب)~$\qty{3.0e16}{m^3}\times917/360 = \qty{76}{m}$، في مقابل \qty{58}{m}.
يستقر معظم جليد غرينلاند على صخر أساسي فوق مستوى سطح البحر، فكلّه
سيضيف ماءً إلى المحيط. ويستقر جزء كبير من جليد أنتاركتيكا على قاع
يقع تحت مستوى سطح البحر: فهذا الجليد يزيح مياه البحر أصلًا، وذوبانه لا يعوّض سوى
الماء الذي يزيحه؛ أمّا الرفوف الجليدية العائمة فلا تضيف شيئًا على الإطلاق.
\end{solution}

\begin{solution}{exo:b3:glaciology:3}
الكسب $\num{6e6}\times1.2 = \qty{7.2e6}{m^3}$، والفقد $\num{4e6}\times2.1 =
\qty{8.4e6}{m^3}$ من الماء: الصافي \qty{-1.2e6}{m^3}، ومتوسط الميزان النوعي
$-1.2\times10^6/10^7 = \qty{-0.12}{m}$ من الماء المكافئ. ونسبة
مساحة التراكم تساوي $6/10 = 0.6$. ويفقد النهر الجليدي كتلة:
فيتقلّص.
\end{solution}

\begin{solution}{exo:b3:glaciology:4}
$1 - 820/917 = \num{0.106}$ و$1 - 830/917 = \num{0.095}$: حوالي عُشر الحجم
ما يزال حيّزًا مساميًّا عندما تنغلق المسام.
\end{solution}

\begin{solution}{exo:b3:glaciology:5}
$\tau_b = 917\times9.81\times300\times\sin\ang{3} = \qty{141.2}{kPa}$. بموجب
مبرهنة الصفيحة، $u_s = \tfrac{2A}{4}\tau_b^3H = 0.5\times\num{1.43e-16}
\times(\num{1.412e5})^3\times300 = \qty{60}{m/yr}$. عند مضاعفة السماكة: $u_s
\propto H^{n+1} = H^4$، أي أسرع ستّ عشرة مرة. وعند نصف الانحدار: $u_s \propto
(\sin\alpha)^3$، أي أبطأ ثماني مرات.
\end{solution}

\begin{solution}{exo:b3:glaciology:6}
عند \qty{-25}{\celsius} (\qty{248.15}{K})، $A = \num{3.61e-13}\exp(-60000/
(8.314\times248.15)) = \qty{8.5e-26}{Pa^{-3}.s^{-1}}$، أو
\qty{2.7e-18}{Pa^{-3}.yr^{-1}}. وعند \qty{-10}{\celsius}، $\num{3.61e-13}
\exp(-27.42) = \qty{4.4e-25}{Pa^{-3}.s^{-1}} = \qty{1.4e-17}{Pa^{-3}.yr^{-1}}$.
عند \qty{0}{\celsius} (الفرع الدافئ) يبلغ \qty{1.4e-16}{Pa^{-3}.yr^{-1}}. وتحت
الإجهادات نفسها تكون السرعات بنسبة عوامل المعدّل:
$\num{1.43e-16}/\num{2.67e-18} = 54$.
\end{solution}

\begin{solution}{exo:b3:glaciology:7}
$h = \sqrt{2\times10^5\times10^6/(917\times9.81)} = \qty{4.7}{km}$، قريب من
\qty{4.9}{km} في أنتاركتيكا: فالنموذج البلاستيكي يعطي سماكة الغطاء الجليدي
انطلاقًا من عرضه وحده.
\end{solution}

\begin{solution}{exo:b3:glaciology:8}
$\tau = 200/5 = \num{40}\text{ سنة}$. ويستمرّ الطول في التكيّف
لعدة عقود، أي بضعة أزمنة استجابة، بعد توقّف الاحترار: لم تلحق الأنهار
الجليدية الحالية بعدُ بمناخ اليوم.
\end{solution}

\begin{solution}{exo:b3:glaciology:9}
(أ)~$h = \sqrt{2\times2.03\times20\times\num{8.64e6}/(917\times\num{3.34e5})}
= \qty{1.51}{m}$. بما أنّ $h^2 \propto t$، فإنّ \qty{3}{m} يستغرق $100\times
(3/1.51)^2 = 393$ يومًا، أي أكثر من شتاء واحد. (ب)~جليد أسمك بكثير من
\qty{3}{m} لا يمكن أن ينمو بالتجمّد خلال شتاء واحد: إنّه جليد نجا
لعدّة شتاءات، أو جليد تكدّس بفعل الرياح والتيارات في الظهرات.
\end{solution}

\begin{solution}{exo:b3:glaciology:10}
(أ)~$H = (1025/917)\times800 = \qty{894}{m}$. (ب)~عند \qty{900}{m} تكون
سماكة \omterm{def:b3:glaciology:grounding-line}{خط الارتكاز} \qty{1006}{m}، أي $9/8$ من السابق؛ ويرتفع الدفق
بمقدار $(9/8)^5 = 1.8$. (ج)~عندما ينحسر الخط على قاع يتعمّق نحو الداخل
فيكون الجليد عند الخط أسمك، فيصرّف أكثر، ويرقّ الجليد
خلفه، فينحسر الخط أكثر: تغذية راجعة موجبة. أمّا الرفّ الجليدي
فيسند الجليد المستقرّ على القاع، فيقلّل الدفق عند \omterm{def:b3:glaciology:grounding-line}{خط الارتكاز} ويمكنه أن
يثبّته في موضعه حتى على منحدر متّجه نحو الأسفل.
\end{solution}

\begin{solution}{exo:b3:glaciology:11}
مع $h(r) = c\sqrt{R - r}$، و$c = \sqrt{2\tau_0/(\rho g)}$، و$V =
\int_0^R2\pi rc\sqrt{R - r}\,\mathrm dr$. بالتعويض $s = R - r$،
\[
\int_0^R(R - s)\sqrt s\,\mathrm ds = \tfrac23R^{5/2} - \tfrac25R^{5/2} =
\tfrac4{15}R^{5/2},
\]
إذن $V = \frac{8\pi}{15}cR^{5/2}$. ومضاعفة $R$
تضرب $V$، وكذلك \omterm{def:b3:glaciology:sle}{مكافئ مستوى سطح البحر}، بمقدار $2^{5/2} = 5.7$.
\end{solution}

\begin{solution}{exo:b3:glaciology:12}
(أ)~مارس: $14.13/16.34 - 1 = \qty{-13.5}{\percent}$؛ سبتمبر: $4.81/7.05
- 1 = \qty{-31.8}{\percent}$. (ب)~في مارس يُمسك حافةَ الجليد إلى حدّ كبير
السواحل المحيطة بالمحيط المتجمّد الشمالي، ويملأ التجمّدُ كل شتاء
الحوض من جديد؛ أمّا الجليد الأرقّ في السنوات الأخيرة فيذوب أكثر
بالكامل كل صيف. وتعمل التغذية الراجعة جليد--ألبيدو في الصيف: فالمياه
المفتوحة تمتصّ نحو \qty{94}{\percent} من ضوء الشمس، والجليد يعكس
\qtyrange{50}{70}{\percent}، فكلّ فقدان لجليد الصيف يُدفئ الماء ويزيد
الذوبان.
\end{solution}

\begin{solution}{pb:b3:glaciology:1}
\textbf{1.} العمود الواقع فوق القاع، ووزنه $\rho gH$ لكل وحدة مساحة، تدفعه
نحو الأسفل المركّبةُ $\rho gH\sin\alpha$، ويوازنها إجهادُ
القصّ على قاعدته؛ وبالنسبة للانحدارات الصغيرة $\sin\alpha \approx \alpha$ (مبرهنة
الصفيحة مع انحدار السطح).

\textbf{2.} $917\times9.81\times580\times0.0190 = \qty{99.1}{kPa}$;
$917\times9.81\times470\times0.0240 = \qty{101.5}{kPa}$; $917\times9.81
\times335\times0.0330 = \qty{99.4}{kPa}$.

\textbf{3.} يبقى الإجهاد ثابتًا ضمن \qty{2}{\percent} بينما
يكاد السمك يتناقص إلى النصف: الغطاء قريب من اللدن، مع $\tau_0 \approx
\qty{100}{kPa}$.

\textbf{4.} $2\times10^5/(917\times9.81) = \qty{22.23}{m}$.

\textbf{5.} $h = \sqrt{22.23\,(R - r)}$: \qty{577}{m}، \qty{472}{m}
و\qty{333}{m} عند 15 و10 و\qty{5}{km} من الحافة، ضمن بضعة أمتار
من قياسات السبر.

\textbf{6.} $\sqrt{22.23\times20000} = \qty{667}{m}$.

\textbf{7.} كما في \cref{exo:b3:glaciology:11}: $V = \int_0^R2\pi r
\sqrt{22.23(R - r)}\,\mathrm dr = \frac{8\pi}{15}\sqrt{2\tau_0/(\rho
g)}\,R^{5/2}$.

\textbf{8.} $\frac{8\pi}{15}\times4.715\times(\num{2e4})^{5/2} =
\qty{4.47e11}{m^3} = \qty{447}{km^3}$.

\textbf{9.} $\pi\times20^2 = \qty{1257}{km^2}$؛ متوسط السمك
$447/1257 = \qty{0.356}{km}$، \qty{356}{m}، أي أكبر قليلًا من نصف
السمك المركزي.

\textbf{10.} $2^{5/2} = 5.7$ مرّة.

\textbf{11.} $A = \num{3.61e-13}\exp(-60000/(8.314\times263.15)) =
\qty{4.45e-25}{Pa^{-3}.s^{-1}}$، مضروبًا في \qty{3.156e7}{s}:
\qty{1.40e-17}{Pa^{-3}.yr^{-1}}.

\textbf{12.} عند \qty{10}{km}، $H = \qty{471.5}{m}$، والانحدار $471.5/20000 =
\num{0.0236}$, $\tau_b = \qty{100}{kPa}$; $u_s = \tfrac{2A}{4}\tau_b^3H =
0.5\times\num{1.40e-17}\times10^{15}\times471.5 = \qty{3.3}{m/yr}$.

\textbf{13.} $q = \tfrac{2A}{5}\tau_b^3H^2 = 0.4\times\num{1.40e-17}
\times10^{15}\times471.5^2 = \qty{1250}{m^2/yr}$؛ متوسط السرعة $q/H =
\qty{2.6}{m/yr}$، أي أربعة أخماس سرعة السطح.

\textbf{14.} $2\pi rq = \pi r^2\bar b$، ومنه $\bar b = 2q/r = 2\times1250/10^4
= \qty{0.25}{m}$ من الجليد سنويًا فوق المنطقة الداخلية ضمن \qty{10}{km}.

\textbf{15.} $h = 450$ حيث $R - r = 450^2/22.23 = \qty{9.11}{km}$، عند
$r = \qty{10.9}{km}$؛ \omterm{def:b3:glaciology:mass-balance}{نطاق التراكم} هو القرص الواقع داخله، $(10.9/
20)^2 = 0.30$ من الغطاء.

\textbf{16.} يرتفع $A$ بمقدار \num{10.2}، وترتفع كذلك السرعات تحت الإجهادات
نفسها. سيتدفّق الجليد نحو الخارج أسرع مما يمدّه به ميزان الكتلة:
فيرقّ الغطاء ويتّسع حتى تعود الإجهادات المنخفضة إلى تطابق ميزان الكتلة
من جديد؛ وللغطاء اللدن الأليَن إجهاد خضوع أصغر
ومقطع أكثر استواءً.

\textbf{17.} $\num{4.47e11}\times917 = \qty{4.10e14}{kg} = \qty{410}{Gt}$.

\textbf{18.} $410/360 = \qty{1.14}{mm}$.

\textbf{19.} $356/0.5 = 711$ سنة.

\textbf{20.} مع انخفاض السطح يصبح الجليد في هواء أدفأ فيذوب أسرع،
والتغذية الراجعة بين الارتفاع--ميزان الكتلة؛ وحين يرقّ الغطاء تتغيّر
الإجهادات والجريان أيضًا. ويتسارع الترقّق.

\textbf{21.} $410/243 = \num{1.7}\text{ سنة}$: الغطاء كله أقلّ من سنتين من
فقد غرينلاند في 2010--2019.

\textbf{22.} يحتفظ الغطاء الجليدي لأركل \omterm{def:b3:glaciology:sle}{بمكافئ مستوى سطح البحر} يبلغ
\textbf{\qty[text-series-to-math=true]{1.1}{mm}} (\qty{1.14}{mm}).
\end{solution}
